\section*{第 \ref{ch:b2:waves-in-media} 章——电磁波在等离子体、导体和电介质中的传播}

\begin{solution}{exo:b2:waves-in-media:1}
$f_p \approx 9\sqrt n$ Hz：\qty{2.8}{MHz}、\qty{9}{MHz}、\qty{9}{GHz}、\qty{2.6e15}{Hz}
($\lambda_p = \qty{115}{nm}$)、\qty{2.8}{kHz}。调幅 \qty{1}{MHz}：仅能穿过
星际空间。调频：穿过电离层，不能穿过鞘层。GPS：
除鞘层外均可通过。光：除铜外均可通过。
\end{solution}

\begin{solution}{exo:b2:waves-in-media:2}
铜：\qty{9.2}{mm}、\qty{2.1}{mm}、\qty{65}{\micro m}、\qty{2.1}{\micro m}；铝在
\qty{1}{MHz} 时为 \qty{84}{\micro m}；海水：\qty{10}{kHz} 时为 \qty{2.3}{m}，
\qty{1}{GHz} 时为 \qty{7}{mm}；土壤：\qty{1}{MHz} 时为 \qty{5}{m}。潜艇位于 \qty{10}{m} 深处时需要 $\delta
\gtrsim \qty{10}{m}$：频率须低于约 \qty{500}{Hz}。
\end{solution}

\begin{solution}{exo:b2:waves-in-media:3}
$n = 1.528$、$1.515$、$1.509$；$v = 1.963$、$1.980$、$\qty{1.987e8}{m/s}$；$n_g = n +
2B/\lambda^2 = 1.545$。折射角 $\arcsin(0.707/n)$：\qty{27.57}{\degree} 与
\qty{27.93}{\degree}：散布为 \qty{0.36}{\degree}。
\end{solution}

\begin{solution}{exo:b2:waves-in-media:4}
$f_p = \qty{4.0}{MHz}$：\qty{5}{MHz} 可通过；\qty{2}{MHz} 被反射，穿透深度为
$c/\sqrt{\omega_p^2 - \omega^2} = \qty{14}{m}$。当 $n = \qty{1.5e12}{m^{-3}}$ 时：\qty{11}{MHz}。密度
随太阳电离通量变化——随小时、季节以及
十一年周期而变，因此最高可用频率也随之改变。
\end{solution}

\begin{solution}{exo:b2:waves-in-media:5}
(a) $1/\gamma\pi a^2 = \qty{5.3}{m\ohm/m}$。(b) $\delta = \qty{65}{\micro m}$；$R \approx 1/\gamma2\pi a\delta
= \qty{41}{m\ohm/m}$，为直流时的 $7.7$ 倍。(c) $\delta = \qty{6.5}{\micro m}$，\qty{0.41}{\ohm/m}，
为 $77$ 倍。(d) 集肤效应占主导后，$R \propto \sqrt f$ 与 $L\omega$ 增长得一样快，
且每匝都会增加电阻；利兹线由许多
彼此绝缘、直径小于 $\delta$ 的细股线绞合而成，因而整根铜线都参与导电。
\end{solution}

\begin{solution}{exo:b2:waves-in-media:6}
$\delta_{\text{Al}} = 11.9$、$0.84$、\qty{0.084}{mm}：$d/\delta = 0.08$、$1.2$、$12$：$0.7$、
$10$、\qty{103}{dB}。在 \qty{50}{Hz} 时该导体是透明的；高
$\mu_r$ 的铁会使磁通发生偏转（磁屏蔽），并使 $\delta$ 缩小为
原来的 $\sqrt{\mu_r}$ 分之一。
\end{solution}

\begin{solution}{exo:b2:waves-in-media:7}
(a) $v_0 = eE_0/m\omega$；$\tfrac12nm\langle v^2\rangle = ne^2E_0^2/4m\omega^2 = \tfrac14\varepsilon_0
E_0^2\,\omega_p^2/\omega^2$。(b) $\tfrac14\varepsilon_0E_0^2$ 与 $\tfrac14\varepsilon_0E_0^2(1 - \omega_p^2/
\omega^2)$。(c) 二者之和为 $\tfrac12\varepsilon_0E_0^2$（括号内为 $1$）；$\langle\Pi
\rangle = E_0^2k/2\mu_0\omega$，两者之比为 $c^2k/\omega = v_g$。(d) $\tfrac14\cdot\tfrac14
/\tfrac12 = 1/8$。
\end{solution}

\begin{solution}{exo:b2:waves-in-media:8}
(a) $1/v_g = (1/c)(1 - \omega_p^2/\omega^2)^{-1/2} \approx (1 + \omega_p^2/2\omega^2)/c$。(b) $\Delta t
= (L\omega_p^2/2c)(1/\omega_1^2 - 1/\omega_2^2)$，其中 $\omega_p^2 = ne^2/\varepsilon_0m$，$\omega =
2\pi f$：常数 $e^2/8\pi^2\varepsilon_0mc = \qty{1.34e-7}{}$（SI）。(c) $nL = \qty{9.3e23}{m^{-2}}$，
$(1/f_1^2 - 1/f_2^2) = \qty{5.7e-18}{s^2}$：$\Delta t = \qty{0.72}{s}$。(d) 柱积分量 $nL$
（色散量），由此可由 $n$ 的模型求得距离。
\end{solution}

\begin{solution}{exo:b2:waves-in-media:9}
(a) $\operatorname{Im}\underline\chi = \omega_p^2\Gamma\omega/[(\omega_0^2 - \omega^2)^2 + \Gamma^2\omega^2]$，其中
$\omega_0^2 - \omega^2 \approx 2\omega_0(\omega_0 - \omega)$：$n'' = \tfrac12\operatorname{Im}\underline\chi$ 即为
所给的洛伦兹线型。(b) $n''(\omega_0) = \omega_p^2/2\omega_0\Gamma$，$\alpha = 2n''\omega_0/c =
\omega_p^2/\Gamma c$。(c) $\omega_p^2 = \qty{3.2e20}{s^{-2}}$：$\alpha = \qty{1.8e4}{m^{-1}}$；
$\ln10/\alpha = \qty{0.13}{mm}$。(d) $n' - 1 \propto (\omega_0 - \omega)/[(\omega_0 - \omega)^2 +
\Gamma^2/4]$：在谱线以下为正，以上为负，在 $\pm\Gamma/2$
谱线两侧范围内递减——该处为反常色散。
\end{solution}

\begin{solution}{exo:b2:waves-in-media:10}
(a) $\underline E = E_0\eu^{-x/\delta}\eu^{\iu(\omega t - x/\delta)}$，$\underline B = \underline k\underline E/\omega
= (1 - \iu)\underline E/\omega\delta$。(b) $\langle\Pi_x\rangle = \tfrac12\operatorname{Re}(\underline E
\underline B^*)/\mu_0 = E_0^2/2\mu_0\omega\delta$（$x = 0$ 处）。(c) $\int_0^\infty\tfrac12\gamma E_0^2
\eu^{-2x/\delta}\dd x = \gamma E_0^2\delta/4 = E_0^2/2\mu_0\omega\delta$，其中用到 $\gamma = 2/\mu_0\omega\delta^2$。
(d) $H_0 = \sqrt2E_0/\mu_0\omega\delta$，因此 $\tfrac12R_sH_0^2 = E_0^2/\gamma\mu_0^2\omega^2\delta^3 =
E_0^2/2\mu_0\omega\delta$。铜在 \qty{1}{GHz} 时：$R_s = 1/\gamma\delta = \qty{8}{m\ohm}$；
$\tfrac12 \times 8 \times 10^{-3} \times 10^4 = \qty{40}{W/m^2}$。
\end{solution}

\begin{solution}{exo:b2:waves-in-media:11}
(a) $\omega_p = \qty{1.4e16}{rad/s}$，$\lambda_p = 2\pi c/\omega_p = \qty{140}{nm}$。(b) $\omega = \qty{3.8e15}{rad/s}
< \omega_p$：$1/\kappa = c/\sqrt{\omega_p^2 - \omega^2} = \qty{23}{nm}$，发生全反射。(c)
仍低于 $\omega_p$：发生反射；仅当波长超过 \qty{140}{nm} 时才能透过（在
钠中为 \qty{210}{nm}）。(d) 碰撞使 $\underline\gamma$ 获得实部，并产生
百分之几的吸收；可见光波段的束缚电子共振
（金、铜）吸收蓝光，使金属呈色。
\end{solution}

\begin{solution}{exo:b2:waves-in-media:12}
(a) $\omega\tau = 0.12$：$\underline\varepsilon_r = 84 - 9.7\iu$；在 \qty{20}{GHz} 时，$\omega\tau = 1$：
$45 - 40\iu$。(b) $\underline n \approx 9.2 - 0.53\iu$；$c/2n''\omega = \qty{1.8}{cm}$：微波炉加热
外层两厘米，其余部分靠热传导加热。(c) $\tfrac12 \times
1.54 \times 10^{10} \times 8.85 \times 10^{-12} \times 9.7 \times 4 \times 10^6 = \qty{2.6}{MW/m^3}$。(d) 在
\qty{20}{GHz} 时吸收长度仅为一毫米：仅表层
会被加热；\qty{2.45}{GHz} 可穿透，且为免许可频段。
\end{solution}

\begin{solution}{pb:b2:waves-in-media:1}
\textbf{1.} $f_p = \qty{9}{MHz}$：GPS 可通过，短波被反射。

\textbf{2.} $k = (\omega/c)\sqrt{1 - \omega_p^2/\omega^2}$：$v_\varphi = \omega/k \approx c(1 + \omega_p^2/
2\omega^2)$，$v_g = \dd\omega/\dd k = c^2k/\omega \approx c(1 - \omega_p^2/2\omega^2)$。

\textbf{3.} $\Delta t = \int(1/v_g - 1/c)\dd z = \int\omega_p^2\dd z/2c\omega^2 = e^2N_T/2\varepsilon_0
mc\omega^2 = (e^2/8\pi^2\varepsilon_0mc)N_T/f^2 = 1.34 \times 10^{-7} \times 2 \times 10^{17}/2.48 \times
10^{18} = \qty{11}{ns}$：\qty{3.2}{m}。

\textbf{4.} $N_T = (\Delta t_1 - \Delta t_2)/K(1/f_1^2 - 1/f_2^2)$；两者之差为
\qty{7}{ns}：$1\%$ 需要 \qty{0.07}{ns}。

\textbf{5.} $v_\varphi$ 超过 $c$ 的量，与 $v_g$ 的不足量相同：载波相位
提前到达，而调制信号滞后到达。
载波相位接收机以相反的符号施加修正。

\textbf{6.} 白天为 \qty{3.2}{m}，夜晚为 \qty{0.3}{m}。

\textbf{7.} $f_p = \qty{0.9}{MHz}$ 考虑碰撞时：中波和短波被
吸收（无线电黑障）；\qty{1.5}{GHz} 下的 GPS 感受到的吸收
$\propto 1/\omega^2$，可忽略。

\textbf{8.} $f_1\Delta t = 1.575 \times 10^9 \times 1.1 \times 10^{-8} = 17$ 个循环。

\textbf{9.} $\delta = \sqrt{2/\mu_0\gamma\omega} = \qty{10}{\micro m}$：一毫米厚的墙相当于
一百个\omterm{def:b2:dispersion-wave-packets:complex}{趋肤深度}。

\textbf{10.} $R_s = 1/\gamma\delta = \qty{0.1}{\ohm}$；$\tfrac12R_sH_0^2 = \qty{45}{W/m^2}$：\qty{45}{W}，
$6\%$。

\textbf{11.} $\eu^{-\pi \times 0.5} = 0.21$：以振幅计 \qty{-14}{dB}，以功率计
\qty{-27}{dB}，且小孔本身的耦合就很弱。

\textbf{12.} 光，$\lambda = \qty{0.5}{\micro m}$，
比小孔小两千倍，故能通过；微波，$\lambda = \qty{12}{cm}$，
比小孔大一百倍。

\textbf{13.} $\underline n = 2 - 0.005\iu$；$c/2n''\omega = \qty{1.9}{m}$：
玻璃几乎不会被加热。

\textbf{14.} $j \approx H_0/\delta = \qty{3e7}{A/m^2}$，$j^2/\gamma = \qty{9e8}{W/m^3}$，
位于十微米厚的薄层内：尖端在数毫秒内升温，发射电子，
使空气电离——形成火花。

\textbf{15.} $\delta = \qty{18}{\micro m}$；截面 $2\pi a\delta = \qty{5.6e-8}{m^2}$，
与 $\pi a^2 = \qty{7.9e-7}{m^2}$ 相比：$R_{\text{交流}} = \qty{0.30}{\ohm/m}$，为直流的十四倍。

\textbf{16.} $RI^2 = \qty{120}{W/m}$；每米的表面积 $\pi d = \qty{3.1e-3}{m^2}$：
$\Delta T \approx \qty{1900}{K}$——会熔化；此类线圈是水冷管。

\textbf{17.} $\delta = \qty{0.14}{mm}$：感应电流与热量，
都集中在十分之一毫米之内；短脉冲加随后的淬火，
使齿轮表面变硬，而芯部仍保持韧性。

\textbf{18.} $R = 1/\gamma2\pi a\delta = \qty{0.10}{\ohm/m}$，
约为原来的三分之一；管内几乎无电流——因此采用管子，
冷却液取代原本会被浪费的铜。

\textbf{19.} $\delta = a$ 时 $f = 1/\pi\mu_0\gamma a^2 = \qty{17}{Hz}$；
低于此频率时，电流均匀分布，直流公式成立。

\textbf{20.} $n^2 \approx 1 + \omega_p^2/\omega_0^2 + (\omega_p^2/\omega_0^2)(\lambda_0/\lambda)^2$：柯西形式，
即 $n(\infty)^2 - 1 = \omega_p^2/\omega_0^2 = 1.25$，对应 $n(\infty) = 1.5$。

\textbf{21.} $n = 1.531$ 和 $1.510$，$v = 1.96$ 和 \qty{1.99e8}{m/s}；$n_g(550) =
1.550$；$n_g(400) - n_g(700) = 0.063$：通过一千米的时间差为 \qty{210}{ns}。

\textbf{22.} $\sin((A + D)/2) = 1.5165 \times 0.5$：$D = \qty{38.6}{\degree}$；$\dd D = 2\sin
(A/2)\dd n/\cos((A + D)/2) = 1.53\,\dd n$：$\Delta D = \qty{1.85}{\degree}$，在 \qty{2}{m} 处为 \qty{6.5}{cm}。

\textbf{23.} 在 \qty{200}{nm} 处频率接近共振，而
$n''$ 增大：不透明。在 \qty{700}{nm} 处的弱共振吸收红光：
绿色玻璃。

\textbf{24.} 蓝光更接近紫外共振：束缚
电子响应更强，$n$ 更大，偏向角也更大。

\textbf{25.} F 层在 \qty{1.5}{GHz}：$\underline\gamma$ 为虚数，\omterm{thm:b2:plane-waves-polarization:wave}{电磁波}传播，
群延迟很小。钢在 \qty{2.45}{GHz}：$\underline\gamma$ 为实数，\omterm{thm:b2:waves-in-media:skin}{趋肤效应}，
发生反射。铜在 \qty{13.56}{MHz}：情况相同，趋肤层厚为
\qty{18}{\micro m}。可见光中的玻璃：$\iu\omega\varepsilon_0\chi$，折射率为实数，
是一种透明的\omterm{def:b2:dispersion-wave-packets:relation}{色散介质}。
\end{solution}
